\begin{portafolioabstract}
	\textbf{Objetivo:} Analizar cómo la educación superior abierta y a distancia desplaza parte del control del aprendizaje hacia el estudiante y convierte la redacción académica en una evidencia central de comprensión, seguimiento y responsabilidad.

	\textbf{Producto elaborado:} Mapa conceptual que organiza, jerarquiza y relaciona conceptos clave del bloque temático mediante nodos, conectores y palabras enlace con función explicativa.

	\textbf{Enfoque de análisis:} Se sostiene que la calidad de la escritura en entornos virtuales no depende solo del manejo formal de la lengua, sino de un sistema compuesto por autogestión, organización del tiempo, responsabilidad académica y criterios de claridad, coherencia y cohesión \citep{nunezcortesEscrituraAcademicaTeoria2015,correaManualBasicoEscritura2018,vitoriaEscrituraAcademicaFormacion2019}.
\end{portafolioabstract}

\subsubsection{Introducción}

En la educación superior abierta y a distancia, escribir bien no es un adorno académico, sino una condición de funcionamiento del aprendizaje. Cuando la mayor parte de la interacción ocurre mediante plataformas, mensajes, consignas y entregables escritos, la redacción deja de ser una habilidad secundaria y se convierte en el medio principal para demostrar comprensión, seguimiento y criterio.

La actividad parte de un problema concreto: explicar por qué la calidad de la escritura en entornos virtuales no depende solo del dominio formal de la lengua, sino de la relación entre autonomía, organización del tiempo, responsabilidad académica y control del propio proceso de estudio. Si esa estructura no se comprende, la escritura se reduce a respuestas fragmentarias, la comunicación pierde eficacia y el aprendizaje deja de ser visible como proceso riguroso \citep{correaManualBasicoEscritura2018,nunezcortesEscrituraAcademicaTeoria2015}.

Desde esta perspectiva, el mapa conceptual funciona como un instrumento de diagnóstico y no solo como un recurso visual. Su utilidad consiste en mostrar que la redacción académica opera como resultado de un sistema: la modalidad virtual exige autogestión; la autogestión necesita disciplina; y ambas condiciones terminan reflejándose en la claridad, coherencia y pertinencia de lo que se comunica.

\subsubsection{Sistema conceptual de la redacción académica en entornos virtuales}

\paragraph{Contexto de análisis}

La educación superior abierta y a distancia redistribuye una parte sustantiva del control del aprendizaje hacia el estudiante. Esa reconfiguración no solo cambia la forma de estudiar; también modifica la forma en que se administra el tiempo, se procesan las instrucciones y se comunica el conocimiento. En este modelo, la escritura funciona como registro operativo del aprendizaje: deja evidencia de comprensión, seguimiento, disciplina y capacidad de respuesta.

Por eso, el problema no se agota en redactar sin errores formales. El punto de fondo consiste en sostener un proceso de estudio capaz de producir textos claros, ordenados y pertinentes dentro de una modalidad flexible, mediada por plataformas y marcada por una menor supervisión inmediata. En otras palabras, la calidad textual no puede separarse de la capacidad de autogestión, porque ambas forman parte del mismo sistema de desempeño académico \citep{vitoriaEscrituraAcademicaFormacion2019,correaManualBasicoEscritura2018}.

\paragraph{Conceptos fundamentales}

\textbf{Educación abierta y a distancia}. Es una modalidad caracterizada por flexibilidad espacial y temporal. Su valor no se encuentra solo en ampliar el acceso, sino en exigir nuevas formas de organización, seguimiento y permanencia dentro del proceso formativo \citep{vitoriaEscrituraAcademicaFormacion2019}.

\textbf{Autogestión del aprendizaje}. Consiste en planificar, supervisar y ajustar el propio estudio. En términos prácticos, funciona como un sistema de control personal: permite administrar tiempos, recursos y estrategias para sostener continuidad y evitar dispersión.

\textbf{Organización del tiempo}. Es la condición operativa que convierte la flexibilidad en avance real. Cuando el tiempo no se estructura, la modalidad virtual pierde eficiencia; cuando se planifica, el aprendizaje adquiere ritmo, secuencia y trazabilidad.

\textbf{Responsabilidad académica}. Implica cumplir compromisos, revisar indicaciones, sostener lectura y cuidar la calidad de lo que se entrega. No es un rasgo abstracto, sino una práctica verificable en la constancia del trabajo y en el nivel de precisión del texto.

\textbf{Aprendizaje autónomo}. Se refiere a la capacidad de avanzar con iniciativa propia, resolver dudas, corregir errores y mantener continuidad sin depender de una vigilancia constante. Esta autonomía no significa aislamiento, sino participación activa y consciente en la propia formación \citep{correaManualBasicoEscritura2018}.

\textbf{Redacción académica}. Es una práctica de comunicación que organiza ideas con orden lógico, lenguaje pertinente y propósito claro. Su función no es adornar el contenido, sino volverlo interpretable, defendible y útil para una comunidad académica \citep{nunezcortesEscrituraAcademicaTeoria2015}.

\textbf{Claridad}. Permite que la idea principal sea comprensible sin ambigüedades innecesarias. Un texto claro reduce ruido, orienta la lectura y facilita la interpretación correcta del mensaje.

\textbf{Coherencia}. Asegura unidad temática y continuidad del argumento. Gracias a ella, cada parte del texto responde a una misma línea de sentido y evita contradicciones internas.

\textbf{Cohesión}. Conecta las partes del discurso mediante relaciones explícitas. Su función es enlazar oraciones, párrafos y conceptos para que el texto avance como sistema y no como fragmentos aislados.

\textbf{Formalidad}. Ajusta el registro al contexto universitario. Esto implica precisión léxica, moderación expresiva y respeto por las convenciones propias de la escritura académica.

\textbf{Comunicación en entornos virtuales}. Es el espacio donde convergen autonomía, responsabilidad y redacción. En esta modalidad, gran parte de la presencia académica se construye a través de mensajes, participaciones y productos escritos; por eso, la calidad de la comunicación depende directamente de la calidad del proceso de estudio.

\paragraph{Síntesis del mapa conceptual}

La actividad original representó estas relaciones mediante un mapa conceptual elaborado en TikZ. Su aporte principal consistió en mostrar que la modalidad virtual impone condiciones, la autogestión traduce esas condiciones en acciones concretas y la redacción académica aparece como la salida verificable de ese proceso. No se trató de una lista de conceptos aislados, sino de una estructura donde cada conector verbal expresaba dependencia, consecuencia o criterio de valoración.

En términos de contenido, el mapa permitió observar tres relaciones decisivas. Primero, que la educación abierta y a distancia exige flexibilidad, pero también demanda organización del tiempo y aprendizaje autónomo. Segundo, que la responsabilidad académica se manifiesta en prácticas concretas como lectura, revisión, corrección y cumplimiento. Tercero, que la calidad textual se valora por criterios como claridad, coherencia, cohesión y formalidad, todos ellos vinculados con la forma en que el estudiante gestiona su propio proceso de estudio.

\clearpage
\begingroup
\pagestyle{empty}
\thispagestyle{empty}
\begin{landscape}
\begin{figure}[p]
\centering
\vspace*{-0.25cm}
\resizebox{0.96\linewidth}{!}{%
\begin{tikzpicture}[
	x=1cm,
	y=1cm,
	concept/.style={
		rectangle,
		rounded corners=3pt,
		draw=dkgray,
		line width=0.45pt,
		fill=lgray,
		align=center,
		text width=2.85cm,
		minimum height=0.85cm,
		inner sep=4pt,
		font=\sffamily\footnotesize
	},
	root/.style={
		concept,
		fill=cardinalred!14,
		draw=cardinalred!70!black,
		text width=4.25cm,
		minimum height=1.25cm,
		font=\bfseries\sffamily\normalsize
	},
	branch/.style={
		concept,
		fill=dkcyan!10,
		draw=dkcyan!60!black,
		text width=2.75cm,
		minimum height=0.95cm,
		font=\bfseries\sffamily\footnotesize
	},
	subbranch/.style={
		rectangle,
		rounded corners=2pt,
		draw=dkgray!75,
		line width=0.35pt,
		fill=white,
		align=center,
		text width=2.55cm,
		minimum height=0.72cm,
		inner sep=3pt,
		font=\sffamily\scriptsize
	},
	edge/.style={
		-{Stealth[length=2.1mm,width=1.5mm]},
		semithick,
		draw=dkgray!90
	},
	relation/.style={
		-{Stealth[length=2.1mm,width=1.5mm]},
		dashed,
		semithick,
		draw=cardinalred!65
	},
	label/.style={
		font=\sffamily\scriptsize,
		text=dkgray,
		fill=white,
		inner sep=1.5pt
	}
]

\node[root] (root) at (0.00,0.00) {Redaccion academica\\en entornos virtuales};

\node[branch] (modalidad) at (-5.82,1.42) {Educacion abierta\\y a distancia};
\node[branch] (autogestion) at (-1.77,3.81) {Autogestion del\\aprendizaje};
\node[branch] (comunicacion) at (1.77,3.82) {Comunicacion en\\entornos virtuales};
\node[branch] (criterios) at (5.84,1.42) {Criterios de\\calidad textual};
\node[branch] (responsabilidad) at (0.00,-3.35) {Responsabilidad\\academica};

\node[subbranch] (modalidad1) at (-8.85,3.15) {Flexibilidad\\espacial y temporal};
\node[subbranch] (modalidad2) at (-9.16,1.43) {Comunicacion\\mediada por texto};
\node[subbranch] (modalidad3) at (-8.85,-0.25) {Menor supervision\\inmediata};

\node[subbranch] (autogestion1) at (-2.35,5.25) {Organizacion\\del tiempo};
\node[subbranch] (autogestion2) at (-4.98,2.94) {Aprendizaje\\autonomo};

\node[subbranch] (comunicacion1) at (2.35,5.25) {Mensajes y\\participaciones};
\node[subbranch] (comunicacion2) at (4.97,2.91) {Productos\\escritos};

\node[subbranch] (criterios1) at (8.85,3.15) {Claridad};
\node[subbranch] (criterios2) at (9.15,1.45) {Coherencia};
\node[subbranch] (criterios3) at (8.85,-0.25) {Cohesion y\\formalidad};

\node[subbranch] (responsabilidad1) at (-3.75,-5.15) {Constancia y\\cumplimiento};
\node[subbranch] (responsabilidad2) at (0.00,-5.45) {Lectura, revision\\y correccion};
\node[subbranch] (responsabilidad3) at (3.75,-5.15) {Texto pertinente\\y riguroso};

\draw[edge] (root.west) -- node[label,above]{se configura en} (modalidad.east);
\draw[edge] (root.north west) -- node[label,left]{requiere} (autogestion.south);
\draw[edge] (root.north east) -- node[label,right]{se expresa en} (comunicacion.south);
\draw[edge] (root.east) -- node[label,above]{se valora por} (criterios.west);
\draw[edge] (root.south) -- node[label,right]{depende de} (responsabilidad.north);

\draw[edge] (modalidad.north west) -- (modalidad1.east);
\draw[edge] (modalidad.west) -- (modalidad2.east);
\draw[edge] (modalidad.south west) -- (modalidad3.east);

\draw[edge] (autogestion.north) -- node[label,left]{implica} (autogestion1.south);
\draw[edge] (autogestion.west) -- node[label,below]{favorece} (autogestion2.east);

\draw[edge] (comunicacion.north) -- (comunicacion1.south);
\draw[edge] (comunicacion.east) -- (comunicacion2.west);

\draw[edge] (criterios.north east) -- node[label,above]{incluyen} (criterios1.west);
\draw[edge] (criterios.east) -- (criterios2.west);
\draw[edge] (criterios.south east) -- (criterios3.west);

\draw[edge] (responsabilidad.south west) -- (responsabilidad1.north);
\draw[edge] (responsabilidad.south) -- (responsabilidad2.north);
\draw[edge] (responsabilidad.south east) -- (responsabilidad3.north);

\draw[relation] (modalidad2.north east) to[out=35,in=180]
	node[label,above,pos=0.45]{favorece}
	(comunicacion1.west);

\draw[relation] (autogestion2.south) to[out=-80,in=150]
	node[label,left,pos=0.45]{fortalece}
	(responsabilidad2.north west);

\draw[relation] (comunicacion2.south east) to[out=-25,in=165]
	node[label,below,pos=0.55]{se ajusta a}
	(criterios2.west);

\end{tikzpicture}
}%
\caption{Mapa conceptual sobre la relacion entre modalidad virtual, autogestion del aprendizaje y calidad de la redaccion academica.}
\label{fig:mapa-conceptual-redaccion-portafolio}
\end{figure}
\end{landscape}
\endgroup
\clearpage

\subsubsection{Conclusión}

El análisis permite afirmar que la redacción académica en entornos virtuales no debe entenderse como una técnica aislada de corrección, sino como el resultado visible de un sistema de trabajo. La modalidad educativa introduce flexibilidad, pero esa flexibilidad solo produce aprendizaje real cuando se convierte en organización, seguimiento, lectura atenta, revisión y responsabilidad sobre lo que se entrega. En ese sentido, escribir bien no es solo presentar un texto limpio; es demostrar control del proceso, capacidad de síntesis y criterio para comunicar conocimiento con precisión.

Desde mi perspectiva, la relación más importante que deja ver esta actividad es la siguiente: la autonomía no vale por sí sola si no produce textos claros, pertinentes y rigurosos. En una carrera como Derecho, donde la argumentación escrita organiza decisiones, interpretaciones y posturas, claridad, coherencia, cohesión y formalidad funcionan como indicadores concretos de disciplina intelectual. Por eso, la redacción académica en ambientes virtuales debe asumirse como una práctica de estructura, responsabilidad y juicio, y no como un requisito accesorio de presentación.
